\scanpage{28}

dans $\unsure{\Sigma|E}$, qu'on peut choisir canoniquement en prenant le quart de grand cercle qui joint $P$ au milieu $2$ du \deleted{segment} l'arc $(\unsure{1|0},\infty)$ de grand cercle \unsure{opposée|opposé} à $0$, ce qui définit un isom du rev universel de $\deleted{\unsure{U}}U$ en $P$, avec celui de $U$ en $\bar P$ sur le quel le groupe précédent opère de façon canonique, donc aussi par transport de structure sur $\widetilde U\deleted{\unsure{(\bar P)}}(P)$, donc il s'envoie dans $\Sigma$.

On trouve que

\[
\dot\sigma_0\mapsto\sigma_0,\quad\dot\tau\longrightarrow\tau_0
\]

est le remontage correspondant. De $=$ pour le pt fixe $-1$ de $\dot\sigma_{\unsure{1}}$, ou $\frac12$ de $\dot\sigma_\infty$, on trouve des scindages partiels sur les ss-groupes de \deleted{\unsure{conj}} stabilité $\Gamma_{-1}\subset\Gamma$ et $\Gamma_2\subset\Gamma$. Il est commode de poser

\[
\deleted{\unsure{P}}\quad 2=0',\quad-1=1',\quad1/2=\infty',
\]

d'où des groupes d'isotropie $\Gamma_{0'}=\{\dot\sigma_0,\tau\}$, $\Gamma_{1'}=\{\dot\sigma_1,\tau\}$, et $\Gamma_{\infty'}=\{\dot\sigma_\infty,\tau\}$, qui se remontent resp en $\sigma_0$ et $\tau_0$, \deleted{\unsure{ou}} $\sigma_1$ et $\tau_1$, en $\sigma_\infty$ et $\tau_\infty$.

Les composantes connexes de $U^\tau$ sont les trois arcs de grand cercle entre $0,1,\infty$, \deleted{qui} \unsure{contient|contenant} chacun un \deleted{pt} et un seul des pts $0',1',\infty'$, et les classes de conjugaison
% archive sheet number: 27
